\chapter{State of the Art}
\label{ch:literature}

\section{Robot Kinematics and Dynamics}
\label{sec:lit_kinematics}

The kinematic and dynamic foundations of serial manipulators are treated
comprehensively in Siciliano et al.\ \cite{Siciliano2009}.
The book covers the Denavit--Hartenberg parameterisation, the geometric
and analytic Jacobian, operational-space dynamics, and redundancy
resolution, and is the standard graduate-level reference in the field.
Craig \cite{Craig2005} provides a complementary introduction that
emphasises homogeneous transformation matrices and forward kinematics;
its frame-convention notation is adopted throughout the present thesis.

Both references are textbooks.
Their treatment of impedance and force control is limited; for
contemporary results, primary literature must be consulted
\cite{Albu2003, Villani2016}.
A formal $SE(3)$ perspective on robot kinematics is given in
Murray et al.\ \cite{Murray1994}.

\section{Impedance Control: Origins and Development}
\label{sec:lit_impedance}

\subsection{Original Formulation}

Hogan \cite{Hogan1985} introduced the impedance control concept by
observing that a robot manipulator interacting with an environment is
best described not by prescribing a rigid trajectory but by specifying
a desired \emph{dynamic relationship} between force and motion.
He formulated this as a target impedance---a virtual mass-spring-damper
at the end-effector---and showed that it unifies position and force
control within a single framework.
This seminal contribution has been cited several thousand times and
remains the theoretical foundation of the present work.

\subsection{Operational-Space Formulation}

Khatib \cite{Khatib1987} developed the operational-space formulation,
in which the full robot dynamics are projected into Cartesian space.
He showed that Cartesian forces can be mapped to joint torques exactly
via the Jacobian transpose---the central equation $\bm{\tau}=\mathbf{J}^\top\mathbf{F}$
used in this thesis---and that the operational-space inertia matrix
can be used to decouple Cartesian axes dynamically.

\textbf{Critical assessment:} While Khatib's full dynamic formulation
yields theoretically superior decoupling, it requires accurate knowledge
of the robot's inertia matrix and introduces computational overhead.
In the present work the simplified quasi-static form is used
($\mathbf{M}=\mathbf{I}$), which is adequate at the low velocities
of the virtual wall task and avoids the sensitivity to model errors
inherent in the dynamic inversion.

\subsection{Admittance Control: An Alternative Approach}

A closely related alternative is \textbf{admittance control}, in which
an outer position/velocity loop receives a force measurement and
generates a modified position reference for an inner position
controller.
Unlike impedance control, admittance control can be layered on top of
any existing position-controlled robot and does not require torque
control capability.
However, it introduces an additional control loop with its own bandwidth
limitation and typically requires an external force/torque sensor.

The Panda's integrated joint torque sensing and FCI torque-control mode
make direct impedance control---as implemented in this thesis---the
more natural and lower-latency choice for this platform.

\subsection{Variable-Impedance and Adaptive Approaches}

Recent research has explored impedance controllers with
\emph{variable} stiffness and damping, adjusted online based on
task requirements or estimated interaction forces.
These approaches improve performance across a wider range of tasks
but introduce stability risks if the parameter update law is not
carefully designed.
The present thesis uses fixed impedance parameters, which guarantees
stability via the Lyapunov argument of \cref{sec:exp_stability} and
is sufficient for the virtual wall task.
Variable-impedance extension is identified as future work in
\cref{sec:future}.

\section{Franka Emika Platform and libfranka}
\label{sec:lit_franka}

The Franka Control Interface \cite{FCI2023} exposes a deterministic
1\,kHz torque-control loop via the open-source \texttt{libfranka}
C++ library.
Unlike ROS-based frameworks, libfranka provides direct access to
joint states, the geometric Jacobian, and model quantities (gravity,
Coriolis, mass matrix) without any middleware overhead.

\textbf{Critical assessment:} The libfranka model is computed from
nominal robot parameters identified by the manufacturer.
For the impedance tasks considered in this thesis (low velocity,
no payload), the model accuracy is sufficient.
Payloads or high-speed motions would require re-identification of
the dynamic parameters.

The FCI imposes a strict 1\,ms callback deadline; violation triggers
an immediate safety halt.
This makes real-time Linux a prerequisite, as documented in
\cref{sec:rt_kernel}.

\section{Summary and Positioning of This Work}
\label{sec:lit_summary}

\cref{tab:lit_comparison} summarises the approaches reviewed above and
positions the present thesis within this landscape.

\begin{table}[htb]
\centering
\caption{Comparison of Cartesian compliance approaches.}
\label{tab:lit_comparison}
\renewcommand{\arraystretch}{1.4}
\begin{tabular}{@{}p{3.2cm}p{2.8cm}p{2.8cm}p{3.5cm}@{}}
\toprule
Approach & Force sensor needed & Requires torque control & Used in this thesis \\
\midrule
Admittance control & Yes & No & No \\
Impedance control (dynamic) & No & Yes & Partial ($\mathbf{M}=\mathbf{I}$) \\
Impedance control (quasi-static) & No & Yes & \textbf{Yes} \\
Variable impedance & Optional & Yes & No (future work) \\
\bottomrule
\end{tabular}
\end{table}

The quasi-static Cartesian impedance controller implemented in this
thesis represents the best trade-off between implementation complexity,
stability guarantees, and suitability for the Panda platform.
The closest related experimental work is Albu-Sch\"affer and Hirzinger
\cite{Albu2003}, who implemented Cartesian impedance control on the
DLR light-weight robot with a similar Jacobian-transpose approach;
their stability analysis corroborates the Lyapunov argument of
\cref{sec:exp_stability}.
A unified treatment of impedance and admittance control is given by
Ott et al.\ \cite{Ott2008}.

%% =============================================================
%%  CHAPTER 3: SYSTEM SETUP AND COMMISSIONING
%% =============================================================
